\subsection{Image statistics and response amplitude}
\label{sec:ampmethods}

\paragraph{Image statistics.} Images were resized and centre-cropped to 224\,px and converted to grey
(Rec.\,709). $L$ is the mean grey value and $C$ its SD (the definitions of Section~\ref{sec:gainmethods}
use Rec.\,601; both correlate almost perfectly). Contrast energy CE is the mean over pixels and 16
filters of the squared response of a zero-mean Gabor bank (4 orientations $\times$ 4 spatial
frequencies, 0.05--0.40 cycles/px, cosine phase, $\sigma = 1/(3 f)$, $41 \times 41$ kernels, reflect
padding). OFF and ON are the mean of $\max(0, \mathrm{ref} - x)$ and of $\max(0, x - \mathrm{ref})$
over pixels, and DD $=$ OFF$/($OFF$+$ON$)$. With the image's own mean as reference OFF equals ON for
every image, so the reference is the mean grey over all images of the data set; DD is then close to a
mirror of $L$ (Spearman \ImgSpearLDD{} in THINGS-720, \ImgSpearLDDNsd{} in NSD).

\paragraph{Response amplitude.} The brain gain $g_i$ is the cosine between pattern $i$ (all voxels of
the ROI, under the given convention) and the mean over the analysed patterns; the model gain $g^M_i$
is the cosine between the raw GAP feature vector of image $i$ and the mean feature vector. With a
common profile left in, two patterns are similar when both have a high amplitude along the profile.
The amplitude RDMs therefore have the sum form $A_{ij} = 2 - g_i - g_j$ and
$A^M_{ij} = 2 - g^M_i - g^M_j$. This form follows from a simulation of the preceding step (step 48,
test (ii): with a common profile $c\,p$ added to amplitude-modulated stimulus patterns, the Spearman
correlation between the correlation-distance RDM and the sum-form RDM rose with $c$). A further
simulation with synthetic patterns, run before the amendment of step 48b and without real data,
showed that the difference form $|g_i - g_j|$ does not capture this structure and that partialling it
does not remove a pure amplitude effect, whereas the sum form does; the difference form is reported
for comparison only (Appendix~\ref{app:conv}). Amplitude dominance is Spearman(RDM, $A$) for the brain
and Spearman(RDM, $A^M$) for a model. The two-sided partial residualises the ranks of the brain RDM on
the ranks of $A$ and the ranks of the model RDM on the ranks of $A^M$ (each with intercept) and
correlates the residuals; the amendment calls this the Spearman of the residuals, and, as in the
partial RSA of step 44 to which it refers, it is computed as the Pearson correlation of the rank
residuals. The moment ratio $m = \lVert \mu \rVert^2 / (\lVert \mu \rVert^2 + \sum_k \sigma_k^2)$ of the
GAP features \citep{evalres} was computed per checkpoint.
